\documentclass[../main]{subfiles}
\begin{document}
\ifSubfilesClassLoaded{\setcounter{page}{4}}{}
\section{Изобретаем «очки-велосипед»}
\label{sec:01_izobretaem_ochki}

Что такое коммунизм? Откуда он взялся, куда зовёт, кого зовёт? Все эти простые вопросы, как показала История, требуют ответов не раз и навсегда, а каждый раз, когда они возникают. Как сказал сценарист Григорий Горин устами «того самого Мюнхгаузена», «любовь – теорема, которую нужно доказывать каждый день». Вот и у нас так же, только коммунизм-то будет теоремой (теорией даже!) посерьёзнее, поглобальнее любви. И надо научиться сперва верно выписывать себе в пункте «дано» для этой теоремы всё подлинно наличное, а не мнимое.

Претензии к самому слову «коммунизм» продиктованы историческим откатом, который произошел с нашим обществом. Вульгаризация и даже превращение в ругательство этого понятия после распада СССР произошла по причине торжества общественных сил, открыто враждебных коммунизму.

До краха Советского Союза уже началось «замыливание» понятия, и потому некоторые теоретики говорили о необходимости разграничения внутри советской экономики и идеологии, проведения линии водораздела, чтобы не принимать капиталистическое инерционное движение за коммунистическую тенденцию. После – по причине того, что общество было ввергнуто в реставрацию капитализма в условиях почти утраченной традиции коммунистической теоретической мысли, когда вопрос стоял о хоть каком-то её сохранении – проблемы коммунизма и движения к нему специально не исследовались. Все взоры были обращены к расцветающему капитализму. И его сторонники, и его противники говорили о нём, как о чём-то реальном, коммунизм же и движение к нему воспринимались как дела прошлого, или, в случае со сторонниками, такого далёкого будущего, до которого ещё дожить нужно, поэтому важно сосредоточиться над тем, чтобы дожить. Коммунизм как теоретическая проблема на десятилетия перестал существовать.

В массовой культуре, далёкой от теоретического мышления, всегда активно воспроизводилось (и воспроизводится сейчас) сознание правящего класса. Этот класс боится коммунизма и ненавидит его, облекая свои страх и ненависть в форму различных социальных мифов, на которых в постсоветских странах растет уже второе поколение. Однако мифологизированные представления о коммунизме не имеют ничего общего с понятием коммунизма. Тут мы находим первую важнейшую точку «враждебности» (которая на самом деле таковой не является). Граждане стран, которые свалились в капитализм не из-за интервенции, а «своими силами», просто не знают уже, что такое коммунизм подлинный, то есть научный, в чём суть идеи коммунизма. И, тем не менее, проблема коммунизма опять становится перед ними как практическая. И та целостная теория, попрание и игнорирование которой сыграло не последнюю роль в сползании назад в капитализм в СССР и других странах соцлагеря, нуждается в возрождении и разъяснении «на пальцах» людям, чьи подлинные интересы она выражает. Речь идёт не только о политически индифферентных тружениках-пролетариях, но и о тех, кто сами себя считают коммунистами или хотя бы левыми, - то есть о политически определившихся как симпатизирующие коммунизму. Последние воспринимают коммунизм как некое «хорошее» общественное устройство, трактуя это «хорошее» скорее на уровне чувственных представлений, часто «от противного» - как отсутствие мерзостей и ужасов современного капиталистического общества. Но положительного ответа на вопрос «что такое коммунизм?» такие его сторонники не могут дать даже себе.

Такое положение дел побуждает именно сейчас создавать эту брошюру как «очки-велосипед», позволяющие снова увидеть коммунизм и понять, что это такое. Что касается противников коммунизма, нужно, как минимум, заставить их выдвигать не глупые, а умные претензии к коммунизму, уже научному.

В год столетия Октября настолько многообразными в публичном пространстве были трактовки роли и поведения большевиков в 1917-м, что мнения, эпитеты, предсказуемо поляризовались: от благородных героев, «спасших государство», до бандитов, это государство (государственность) ненавидевших и уничтоживших. Однако главное при такой поляризации ускользнуло: большевики и, конечно же, Ленин были мыслителями и в своих действиях руководствовались не только желаниями, предпочтениями и т.д., но и теорией, ставшей итогом всего предшествующего развития человеческой мысли. Враги социализма и коммунизма любят называть Октябрь и дальнейшее социалистическое строительство экспериментом, поставленным большевиками над народами… Так вот эти враги – гораздо ближе к истине, чем далёкие от теории друзья.

Именно экспериментом, если конечно не понимать это слово глупо, как внешние манипуляции, была Великая Октябрьская Революция. Ничем другим в силу своей беспрецедентности она и не могла быть. Экспериментом, проведённым по всем научным традициям – ведь не бывает экспериментов без гипотез. И такая гипотеза, разработанная Лениным и его сторонниками и основанная на наследии Маркса и Энгельса, была - Мировая Революция. Начавшись в России, пролетарская революция должна была всколыхнуть индустриально развитые европейские страны, дав толчок для преобразования их во Всемирную республику Советов. Далеко не все научные гипотезы подтверждаются экспериментом в исходной формулировке (несмотря на то, что идея актуальности пролетарских революций в общем остаётся верной и сейчас), однако и это экспериментальное движение по намеченному даже отчасти неподтверждённой гипотезой пути продвигает науку (социализм) вперёд.

Предлагаемая метафора лишь возвращает нас к иному способу видения революции, военного коммунизма, нэпа, индустриализации и так далее, призывает рассматривать современное общество и тенденции его развития через призму коммунизма как важнейшей науки в системе всех общественных наук, необходимых для того, чтобы можно было действовать разумно в современном мире. Чтобы революционный класс мог преодолеть подчинение постоянной эксплуатации, преодолеть в себе угнетённый класс капиталистического общества и встать на собственную классовую позицию, осознать и разумно изменить общественные отношения в сторону уничтожения всякого угнетения и эксплуатации человека человеком, в сторону коммунизма.

Вот почему высмеянные в другом контексте Маяковским «очки-велосипед» теории коммунизма нам помогут только в том случае, если мы наберёмся усидчивости «изобрести» их заново (воспроизвести): пройти Марксов путь к Коммуне и Интернационалу, ленинский теоретический путь к Октябрю, а затем и весь советский путь, не меняя «средств передвижения» - революционной теории.

Заново – значит сделать своим наследие науки коммунизма в современных, доставшихся нам, исторических условиях.
\end{document}
